% Local-only guide; distinguish executable baseline from planned scope.
\chapter{INSTALLATION AND USER GUIDE}\label{app:installation}

\section{Installation Requirements}
\subsection{Hardware Requirements}
CPU inference runs on one local computer. Record actual OS, CPU, RAM and free disk in the local test
report; batch/import/project storage requirements remain to be measured. A GPU is not required by the
runtime application. Offline research training has separate requirements.
\subsection{Software Requirements}
The baseline uses a Python virtual environment and pinned backend dependencies, Node/npm with the
frontend lockfile, and locally prepared reference data/model weights. The target adds local PostgreSQL
and a job worker. Exact versions are recorded by the implementation's requirements/lockfiles;
PostgreSQL/job integration is still pending. The research runner additionally uses its own packages
and MMseqs2; these are not all required for serving predictions.

\section{Local System Execution}
There is one local web application. Public hosting, VPS, domain and TLS deployment are outside the
required workflow. All application listeners must be bound to loopback.

\subsection{Existing Baseline Commands}
From the repository root, create/install the baseline backend environment with
\nolinkurl{make -C topics/degron-footprint/system/app venv}, then run
\nolinkurl{make -C topics/degron-footprint/system/app dev-api}.
The API listens on \nolinkurl{http://127.0.0.1:8000}. On first run this command prepares a SQLite
development database from the local research checkout; it does not start the planned PostgreSQL
workspace or durable worker. The existing development ETL's default ClinVar subset must be recorded
when interpreting annotations.

In a second terminal enter \nolinkurl{topics/degron-footprint/system/app/frontend}, run
\nolinkurl{npm ci}, then \nolinkurl{npm run dev -- --host 127.0.0.1}. Open
\nolinkurl{http://127.0.0.1:5173}. Use the real backend; mock mode contains synthetic scores.
These commands are read from the baseline guides and were not executed as application acceptance
in this document-update task.

\subsection{Target Local Workflow}
\textbf{[PENDING: implemented migrations and startup commands for the local PostgreSQL workspace,
worker and new screens.]}
The final guide must reproduce reference preparation, local startup, a saved project, a peptide/protein
job, a committed observation dataset, comparison/export, shutdown and reopening on the same machine.
Copy database and workspace files together for recovery. Docker may package these same components;
it is not a separate public/internal edition.

\section{User Guide}
The target sequence is: create/open a project; choose single peptide, peptide batch or whole protein;
review syntax/domain/coverage; submit and inspect per-tile results; view ClinVar/reference evidence;
import CSV/TSV with assay/units/replicates and fix validation errors; explicitly commit; compare compatible
observations; export scores and provenance; reopen history. Unsaved single-tile previews are transient;
saved work persists until deletion. Deletion presents dependencies and requires confirmation.
\textbf{[PENDING: screenshots and an executed local walkthrough after implementation.]}
